% Shared section reading content fragment. Do not compile this file directly.
\section{Curriculum Overview and Subject Matter Scope}

This document establishes the official academic reading material and syllabus scope for \textbf{CS 201 --- Algorithms II - Design of Algorithms}. Learners will engage with theoretical foundations, mathematical derivations, practical methodologies, and production-grade applications across the entire sequence.

\begin{InsightBox}{Official Syllabus Scope}
  \textbf{Curriculum Topics}: Advanced algorithm design paradigms: dynamic programming, greedy strategies, network flow, amortized analysis, string matching, and NP-completeness. Direct continuation of Algorithms I and prerequisite for Algorithms III.
\end{InsightBox}

\subsection{Instructional Methodology and Reading Expectations}

Each chapter in this reference text is structured to provide rigorous treatment of both principles and applications. Learners are expected to:
\begin{enumerate}
  \item Master foundational definitions, theorems, and mathematical models presented in each section.
  \item Trace through worked numerical and analytical examples in detail.
  \item Implement core algorithms and computational workflows using standard open-source tools.
  \item Complete end-of-chapter exercises and review analytical edge cases.
\end{enumerate}

\ContentPendingChapter{This reading orientation establishes the approved course syllabus scope. Comprehensive narrative expositions, proofs, worked engineering case studies, and assessment problem sets will be progressively integrated by the course editorial team.}
